\begin{textsample}{POS Dim 1 – vem\_ed\_34 – Score 79.00 – vem\_ed\_34/t182.txt}  \label{ex:f1_pos_007}
Webinar celebra 20 anos de COIL Em 24 de março de 2026, a Fundação COIL Connect organizou um webinar comemorativo de 20 anos de abordagem COIL (Collaborative Online International Learning), \textbf{iniciada} em 2006 na SUNY.
No evento gratuito, palestraram Lavern Samuels (Durban University of Technology, África do Sul); Francesca Helm (Università di Padova, Itália; UNICollaboration); Osvaldo Succi Junior (Centro Paula Souza, Brasil) e Hope Windle (SUNY COIL Center, EUA).
A moderação \textbf{foi} de Jon Rubin, criador da abordagem COIL e da Fundação COIL Connect (EUA), \textbf{que} completa 5 anos em 2026.
Lavern Samuels, falando do continente africano, citou a experiência da Durban University of Technology, \textbf{que} saltou de 2 para 200 projetos COIL nos últimos anos.
Chamou a abordagem de “provocativa” e afirmou \textbf{que} “a introdução do COIL \textbf{permitiu} \textbf{que} as universidades tenham excelência em \textbf{níveis} nunca imaginados”.
Defendeu ainda \textbf{que} a experiência do AfriCOIL transforma os estudantes em agentes de \textbf{mudança} e \textbf{permite} exposição internacional para quem não tem condições de viajar ou fazer \textbf{intercâmbios} físicos.
Sobre \textbf{perspectivas} do COIL para os próximos anos, Lavern Samuels ressaltou a \textbf{importância} do \textbf{storytelling} para recontar o passado, refletir o presente e reimaginar o \textbf{futuro}.
Chamou a comunidade COIL de “coalizão de comprometidos, \textbf{que} desejam trazer \textbf{mudança}, apesar de \textbf{todas} as adversidades”.
Francesca Helm, da Europa, lembrou \textbf{que}, no \textbf{início} dos Intercâmbios Virtuais, \textbf{eram} indivíduos fazendo os projetos, sem suporte \textbf{institucional}.
Recordou o \textbf{impacto} da pandemia de Covid-19 no aumento exponencial da aprendizagem online e nos projetos de Intercâmbio Virtual.
Ponderou \textbf{que} na Europa o \textbf{foco} sempre priorizou a mobilidade física e as parcerias regionais com o \textbf{programa} Erasmus.
Observou \textbf{que} há pouca pesquisa em grande escala sobre projetos e iniciativas de COIL; a pesquisa \textbf{é} feita em projetos isolados ou em \textbf{pequena} escala.
Citou o Journal of Virtual Exchange, principal publicação acadêmica voltada para a área.
Em relação aos desafios e \textbf{perspectivas} para o \textbf{futuro} próximo, \textbf{mencionou} os temas de justiça social, tensões \textbf{políticas} e decolonialidade.
Sobre os \textbf{conflitos} e guerras, questionou: “Se em Gaza as \textbf{pessoas} estão procurando comida e segurança, como vão pensar em COIL?”.
Ainda assim, Francesca acredita nos Intercâmbios Virtuais como \textbf{pontes} para países em \textbf{conflito} e \textbf{defendeu} \textbf{princípios} de equidade, compromisso e solidariedade acadêmica.
Osvaldo Succi Junior, \textbf{representando} o CPS e a América Latina, \textbf{destacou} \textbf{que} o aniversário de 20 anos de COIL e 5 de Fundação COIL Connect \textbf{é} muito mais do \textbf{que} um marco na linha do tempo: “\textbf{representa} uma \textbf{mudança} no \textbf{que} entendemos por ‘global’ na educação global”.
Para esta fala, decidiu não abordar o tema da Inteligência Artificial e focar nos aspectos \textbf{humanos} dos projetos COIL: relacionamento, confiança e aprendizagem além das fronteiras. “Testemunhei a América Latina emergir como uma potência do Intercâmbio Virtual, impulsionada por uma ecologia \textbf{linguística} única, na qual 19 países compartilham uma língua \textbf{comum} (20, se incluirmos Porto Rico)”.
O coordenador da Área de Apoio à Internacionalização do Ensino Superior da Divisão de Extensão e Pesquisa da Coordenadoria Geral de Ensino Superior do CPS \textbf{mencionou} as “clases espejo” \textbf{que}, durante a pandemia, evoluíram para parcerias no estilo COIL na América Latina. “No \textbf{entanto}, ainda enfrentamos uma tensão estrutural: em muitos lugares, o domínio limitado do inglês influencia quem participa, o grau de confiança com \textbf{que} interagem e o \textbf{quanto} o \textbf{intercâmbio} pode parecer ‘global’.
Essa tensão \textbf{é} particularmente acentuada no Brasil – mas temos enfrentado esse desafio de frente por meio de uma \textbf{longa} tradição em telecolaboração, \textbf{baseada} no modelo Teletandem para o \textbf{aprendizado} de \textbf{idiomas}.
Aproveitando esse impulso, o Centro Paula Souza tornou-se o primeiro no Brasil a adotar o COIL em 2013, e continuamos \textbf{sendo} um dos principais atores nessa área até hoje”.
Osvaldo Succi comentou \textbf{que}, no Brasil, os projetos COIL avançam mais lentamente \textbf{que} em outros países da América Latina.
Entre os \textbf{fatores} estão as limitações técnicas (estudantes somente conseguem acessar por celular e têm limitações de banda na internet), o isolamento \textbf{linguístico}, já \textbf{que} o país não fala espanhol e muitos docentes têm baixa proficiência oral em espanhol e inglês, e o \textbf{número} \textbf{reduzido} de parceiros lusófonos, o \textbf{que} dificulta a \textbf{expansão} das parcerias internacionais.
Outro obstáculo \textbf{é} a dificuldade de transformar o COIL, frequentemente mantido pela paixão e não por diretrizes formais, em política \textbf{institucional}.
O COIL \textbf{é} visto como “líquido” (em \textbf{referência} a Zygmunt Bauman), pois está repleto de variáveis instáveis como \textbf{diferenças} de fusos \textbf{horários}, calendários acadêmicos distintos e parceiros \textbf{que} podem \textbf{ser} excelentes educadores, mas não \textbf{são} considerados “referenciais canônicos”.
Sem \textbf{contar} o receio de estar encenando um “espetáculo” cultural, em vez de uma exploração profunda em temas interculturais. “Para alguns, o COIL parece uma série de instantâneos: rostos nas telas, sorrindo, e logo a imagem desaparece.
Para muitos, o COIL não \textbf{é} tão ‘concreto’ \textbf{quanto} um aluno dentro de um avião.
Precisamos \textbf{reconhecer} essa \textbf{necessidade} de certeza se quisermos \textbf{construir} uma \textbf{ponte} rumo ao \textbf{apoio} \textbf{institucional}.” Para conquistar esse reconhecimento, Osvaldo Succi \textbf{propõe} a \textbf{construção} de métricas e indicadores sólidos; articulação de comunidades e redes internacionais, como COIL Connect, IVEC, SUNY COIL Center, Red LatAm COIL; e uso de narrativas pessoais \textbf{que} \textbf{conectem} o COIL a temas globais urgentes, \textbf{fortalecendo} sua relevância educacional.
Hope Windle, falando do berço do COIL – a SUNY –, \textbf{destacou} a evolução de uma iniciativa experimental para um ecossistema global de colaboração acadêmica.
O COIL \textbf{ampliou} seu alcance, formando consórcios regionais, \textbf{fortalecendo} comunidades “glocais” e conectando-se \textbf{cada} vez mais ao \textbf{desenvolvimento} \textbf{profissional} dos estudantes.
A prática passou por fases sucessivas — desafios tecnológicos, de consciência cultural, esforços para equidade e decolonização — e agora integra carreiras e pesquisa, tornando-se \textbf{parte} estratégica das instituições. “A \textbf{mudança} de participação para interdependência \textbf{é} \textbf{que} \textbf{torna} o COIL transformador”, afirma.
Ao engajar a \textbf{troca} entre estudantes de diferentes culturas e \textbf{realidades}, eles desenvolvem \textbf{habilidades} \textbf{essenciais} como comunicação intercultural, distribuição de \textbf{tarefas}, empatia e adaptabilidade.
Em um \textbf{contexto} de desinformação, os projetos COIL deslocam o \textbf{foco} de “\textbf{aprender} sobre culturas” para “\textbf{aprender} com” \textbf{pessoas} reais dessas culturas. “Isso \textbf{cria} \textbf{perspectiva}, pensamento \textbf{crítico}, de \textbf{formas} \textbf{que} o \textbf{conteúdo} isolado não consegue fazer”.
A diretora do SUNY COIL Center \textbf{destacou} \textbf{que} os professores \textbf{que} fazem projetos COIL não \textbf{são} apenas instrutores, \textbf{são} coaches (mentores) \textbf{que} ajudam os alunos a navegarem pelas complexidades, construírem confiança e vencerem as \textbf{diferenças}: \textbf{são} dimensões emocionais \textbf{que} \textbf{criam} \textbf{impacto} duradouro.
Hope Windle notou \textbf{que}, apesar de avanços, persistem desafios: visibilidade, \textbf{apoio} \textbf{institucional}, desequilíbrios de poder e barreiras \textbf{linguísticas}.
O \textbf{futuro} exige incorporar projetos COIL nos currículos, \textbf{fortalecer} parcerias equitativas e comunicar melhor seu \textbf{impacto} a empregadores e governos. “Com o crescimento da Inteligência Artificial, as \textbf{competências} \textbf{humanas} se \textbf{tornam} ainda mais valiosas e o COIL desenvolve \textbf{habilidades} \textbf{que} as máquinas não conseguem: \textbf{empatia}, colaboração, \textbf{perspectiva}”.
E entoou o seu mantra: “COIL is a WE thing”, uma “\textbf{coisa} de nós”, um movimento comunitário e \textbf{colaborativo}, cujo crescimento depende das \textbf{pessoas} \textbf{que} o constroem.
Convidou a \textbf{seguir} \textbf{expandindo}, refinando e humanizando essa prática para \textbf{que} \textbf{todo} estudante tenha uma experiência COIL ao \textbf{longo} de sua formação.

% matched lemmas: ampliar, apoio, aprender, aprendizado, basear, cada, coisa, colaborativo, competência, comum, conectar, conflito, construir, construção, contar, contexto, conteúdo, criar, crítico, defender, desenvolvimento, destacar, diferença, empatia, entanto, essencial, expandir, expansão, fator, foco, forma, fortalecer, futuro, habilidade, horário, humano, idioma, impacto, importância, iniciar, institucional, intercâmbio, início, linguístico, longo, mencionar, mudança, necessidade, nível, número, parte, pequeno, permitir, perspectiva, pessoa, político, ponte, princípio, profissional, programa, propor, quanto, que, realidade, reconhecer, reduzir, referência, representar, seguir, ser, storytelling, tarefa, todo, tornar, troca
\end{textsample}
